% Vietnamese translation for OI Public Library
% Source-PDF: USACO/USACO全部中英文译题83P.pdf
\documentclass[11pt]{article}
\usepackage{oipl-vn}

\title{USACO Training Problem List}
\author{}
\date{}

\newcommand{\usacoproblem}[3]{%
  \subsection*{#1 \hfill {\normalsize\texttt{#2}}}%
  \origtitle{#3}%
}
\newcommand{\ioheading}[1]{\paragraph{#1}}

\begin{document}
\maketitle

\origtitle{USACO Train Problem List, collected by BirDOR}
Tài liệu này chuyển các phát biểu bài trong PDF nguồn sang tiếng Việt, giữ nguyên tên bài và mã chương trình chính thức của USACO. Các ví dụ vào/ra được đặt ở phụ lục theo đúng thứ tự xuất hiện trong PDF để tránh làm đứt mạch danh sách bài.

\section*{Chương 1}
\subsection*{Section 1.1}
\usacoproblem{Your Ride Is Here}{ride}{Your Ride Is Here}
\ioheading{Nhiệm vụ.}
Mỗi tên sao chổi và tên nhóm được đổi thành một số bằng tích các giá trị chữ cái, với \texttt{A}=1 và \texttt{Z}=26. Nếu hai tích có cùng phần dư khi chia cho 47 thì nhóm được chiếc UFO đi sau sao chổi đó đón đi.

\ioheading{Dữ liệu vào.}
Vào gồm tên sao chổi và tên nhóm, mỗi tên là chuỗi chữ hoa dài từ 1 đến 6, không chứa khoảng trắng hay dấu câu.

\ioheading{Kết quả.}
In \texttt{GO} nếu hai tên khớp theo quy tắc trên, ngược lại in \texttt{STAY}.

\usacoproblem{Greedy Gift Givers}{gift1}{Greedy Gift Givers}
\ioheading{Nhiệm vụ.}
Một nhóm bạn tặng quà cho nhau. Mỗi người có một khoản tiền và chia đều nhiều nhất có thể cho danh sách người nhận của mình; phần không chia hết được giữ lại. Cần tính số tiền ròng mỗi người nhận được trừ đi số tiền đã cho.

\ioheading{Dữ liệu vào.}
Dòng đầu là số người \texttt{NP}. Sau đó là \texttt{NP} tên theo thứ tự cần in. Với mỗi người tặng, dữ liệu cho tên người tặng, số tiền, số người nhận, rồi danh sách người nhận.

\ioheading{Kết quả.}
In mỗi tên theo thứ tự ban đầu, kèm số tiền ròng của người đó.

\usacoproblem{Friday the Thirteenth}{friday}{Friday the Thirteenth}
\ioheading{Nhiệm vụ.}
Đếm trong \texttt{N} năm, từ ngày 1/1/1900 đến 31/12/(1900+\texttt{N}-1), ngày 13 của mỗi tháng rơi vào từng thứ trong tuần bao nhiêu lần. Cần dùng đúng quy tắc năm nhuận Gregorian: năm chia hết cho 4 là nhuận, trừ năm thế kỷ không chia hết cho 400.

\ioheading{Dữ liệu vào.}
Một số nguyên \texttt{N}, với \texttt{N} không quá 400.

\ioheading{Kết quả.}
In bảy số theo thứ tự thứ Bảy, Chủ nhật, thứ Hai, ..., thứ Sáu.

\usacoproblem{Broken Necklace}{beads}{Broken Necklace}
\ioheading{Nhiệm vụ.}
Cho vòng cổ gồm \texttt{N} hạt màu đỏ, trắng, xanh. Sau khi bẻ vòng ở một vị trí, có thể nhặt liên tiếp sang trái và sang phải; hạt trắng có thể đóng vai trò đỏ hoặc xanh. Hỏi nhiều nhất nhặt được bao nhiêu hạt.

\ioheading{Dữ liệu vào.}
Dòng đầu là \texttt{N} (\texttt{3 <= N <= 350}); dòng sau là chuỗi ký tự \texttt{r}, \texttt{w}, \texttt{b}.

\ioheading{Kết quả.}
In số hạt tối đa có thể thu được, không vượt quá \texttt{N}.

\subsection*{Section 1.2}
\usacoproblem{Milking Cows}{milk2}{Milking Cows}
\ioheading{Nhiệm vụ.}
Có nhiều nông dân vắt sữa trong các khoảng thời gian. Cần tìm khoảng thời gian dài nhất luôn có ít nhất một người vắt sữa và khoảng nghỉ dài nhất không có ai vắt sữa giữa lúc bắt đầu và kết thúc toàn bộ công việc.

\ioheading{Dữ liệu vào.}
Dòng đầu là số khoảng \texttt{N}. Mỗi dòng sau chứa thời điểm bắt đầu và kết thúc của một khoảng.

\ioheading{Kết quả.}
In hai số: thời gian vắt sữa liên tục dài nhất và thời gian nhàn rỗi dài nhất.

\usacoproblem{Transformations}{transform}{Transformations}
\ioheading{Nhiệm vụ.}
Cho một mẫu hình vuông ban đầu và mẫu sau biến đổi. Xác định phép biến đổi đầu tiên trong danh sách: quay 90 độ, 180 độ, 270 độ, phản chiếu ngang, phản chiếu rồi quay, không đổi, hoặc không hợp lệ.

\ioheading{Dữ liệu vào.}
Dòng đầu là \texttt{N}; tiếp theo là \texttt{N} dòng mẫu ban đầu và \texttt{N} dòng mẫu đích.

\ioheading{Kết quả.}
In số hiệu từ 1 đến 7 ứng với phép biến đổi đúng theo thứ tự ưu tiên của đề.

\usacoproblem{Name That Number}{namenum}{Name That Number}
\ioheading{Nhiệm vụ.}
Một số trên bàn phím điện thoại có thể tương ứng với các tên trong từ điển, với các chữ cái \texttt{Q} và \texttt{Z} không được dùng. Cần liệt kê mọi tên đổi ra đúng dãy chữ số đã cho.

\ioheading{Dữ liệu vào.}
Một dòng chứa dãy chữ số. Từ điển \texttt{dict.txt} là dữ liệu phụ chuẩn của bài.

\ioheading{Kết quả.}
In các tên hợp lệ theo thứ tự từ điển; nếu không có tên nào thì in \texttt{NONE}.

\usacoproblem{Palindromic Squares}{palsquare}{Palindromic Squares}
\ioheading{Nhiệm vụ.}
Với cơ số \texttt{B}, tìm các số từ 1 đến 300 sao cho bình phương của chúng, viết trong cơ số \texttt{B}, là palindrome.

\ioheading{Dữ liệu vào.}
Một số nguyên \texttt{B}.

\ioheading{Kết quả.}
Mỗi dòng in số ban đầu và bình phương của nó trong cơ số \texttt{B}.

\usacoproblem{Dual Palindromes}{dualpal}{Dual Palindromes}
\ioheading{Nhiệm vụ.}
Tìm \texttt{N} số nguyên lớn hơn \texttt{S} mà biểu diễn của chúng là palindrome trong ít nhất hai cơ số từ 2 đến 10.

\ioheading{Dữ liệu vào.}
Hai số \texttt{N} và \texttt{S}.

\ioheading{Kết quả.}
In \texttt{N} số tìm được theo thứ tự tăng dần, mỗi số một dòng.

\subsection*{Section 1.3}
\usacoproblem{Mixing Milk}{milk}{Mixing Milk}
\ioheading{Nhiệm vụ.}
Farmer John cần mua một lượng sữa nhất định từ nhiều nông dân, mỗi người có giá bán và lượng sữa tối đa. Hãy mua đủ sữa với chi phí nhỏ nhất.

\ioheading{Dữ liệu vào.}
Dòng đầu chứa lượng cần mua và số nông dân. Mỗi dòng sau chứa giá mỗi đơn vị và lượng sữa của một nông dân.

\ioheading{Kết quả.}
In tổng chi phí nhỏ nhất.

\usacoproblem{Barn Repair}{barn1}{Barn Repair}
\ioheading{Nhiệm vụ.}
Có các chuồng bò đánh số liên tiếp và chỉ một số chuồng có bò. Dùng nhiều nhất \texttt{M} tấm ván để che tất cả chuồng có bò, sao cho tổng số chuồng bị che là nhỏ nhất.

\ioheading{Dữ liệu vào.}
Dòng đầu chứa \texttt{M}, tổng số chuồng và số chuồng có bò; các dòng sau là vị trí chuồng có bò.

\ioheading{Kết quả.}
In tổng chiều dài nhỏ nhất của các tấm ván.

\usacoproblem{Calf Flac}{calfflac}{Calf Flac}
\ioheading{Nhiệm vụ.}
Trong một văn bản, tìm palindrome dài nhất nếu bỏ qua ký tự không phải chữ cái và không phân biệt hoa thường. Khi in kết quả phải giữ nguyên đoạn gốc.

\ioheading{Dữ liệu vào.}
Toàn bộ tệp vào là văn bản, có thể chứa xuống dòng và dấu câu.

\ioheading{Kết quả.}
In độ dài palindrome theo số chữ cái, rồi in đoạn văn bản gốc tương ứng.

\usacoproblem{Prime Cryptarithm}{crypt1}{Prime Cryptarithm}
\ioheading{Nhiệm vụ.}
Đếm các phép nhân dạng ba chữ số nhân hai chữ số sao cho mọi chữ số xuất hiện trong hai thừa số, hai tích riêng và tích cuối đều thuộc tập chữ số cho trước; hai tích riêng phải có ba chữ số và tích cuối phải có bốn chữ số.

\ioheading{Dữ liệu vào.}
Dòng đầu là số chữ số cho phép, dòng sau là tập chữ số.

\ioheading{Kết quả.}
In số phép nhân hợp lệ.

\subsection*{Section 1.4}
\usacoproblem{The Clocks}{clocks}{The Clocks}
\ioheading{Nhiệm vụ.}
Có 9 đồng hồ chỉ 12, 3, 6 hoặc 9 giờ. Chín thao tác cho trước sẽ xoay các nhóm đồng hồ thêm 3 giờ. Cần tìm dãy thao tác ngắn nhất theo thứ tự từ điển đưa tất cả đồng hồ về 12 giờ.

\ioheading{Dữ liệu vào.}
Ba dòng, mỗi dòng ba giá trị giờ hiện tại.

\ioheading{Kết quả.}
In dãy số thao tác, cách nhau bởi dấu cách.

\usacoproblem{Packing Rectangles}{packrec}{Packing Rectangles}
\ioheading{Nhiệm vụ.}
Cho bốn hình chữ nhật, mỗi hình có thể xoay. Hãy ghép chúng không chồng nhau thành một hình chữ nhật bao ngoài có diện tích nhỏ nhất, rồi liệt kê mọi kích thước đạt diện tích đó.

\ioheading{Dữ liệu vào.}
Bốn dòng, mỗi dòng chứa chiều rộng và chiều cao của một hình chữ nhật.

\ioheading{Kết quả.}
In diện tích nhỏ nhất, sau đó là các cặp kích thước theo thứ tự tăng, luôn ghi cạnh nhỏ trước.

\usacoproblem{Arithmetic Progressions}{ariprog}{Arithmetic Progressions}
\ioheading{Nhiệm vụ.}
Một bisquare là số dạng \(p^2+q^2\) với \(0\le p,q\le M\). Cần tìm mọi cấp số cộng dài \texttt{N} chỉ gồm các bisquare.

\ioheading{Dữ liệu vào.}
Hai dòng chứa \texttt{N} và \texttt{M}.

\ioheading{Kết quả.}
In các cặp điểm đầu và công sai theo công sai tăng dần rồi điểm đầu tăng dần; nếu không có thì in \texttt{NONE}.

\usacoproblem{Mother's Milk}{milk3}{Mother's Milk}
\ioheading{Nhiệm vụ.}
Ba thùng có dung tích cố định, ban đầu chỉ thùng C đầy sữa. Có thể rót từ thùng này sang thùng khác cho tới khi nguồn cạn hoặc đích đầy. Hỏi khi thùng A rỗng thì thùng C có thể chứa những lượng nào.

\ioheading{Dữ liệu vào.}
Một dòng chứa dung tích ba thùng A, B, C.

\ioheading{Kết quả.}
In các lượng sữa có thể có trong thùng C theo thứ tự tăng dần.

\subsection*{Section 1.5}
\usacoproblem{Number Triangles}{numtri}{Number Triangles}
\ioheading{Nhiệm vụ.}
Cho tam giác số. Bắt đầu từ đỉnh, mỗi bước đi xuống một trong hai số kề ngay dưới; cần tổng lớn nhất trên một đường đi tới đáy.

\ioheading{Dữ liệu vào.}
Dòng đầu là số hàng, sau đó là các hàng của tam giác.

\ioheading{Kết quả.}
In tổng lớn nhất.

\usacoproblem{Prime Palindromes}{pprime}{Prime Palindromes}
\ioheading{Nhiệm vụ.}
Liệt kê mọi số vừa là số nguyên tố vừa là palindrome trong đoạn \([a,b]\).

\ioheading{Dữ liệu vào.}
Một dòng chứa hai số \texttt{a} và \texttt{b}.

\ioheading{Kết quả.}
In mỗi số hợp lệ trên một dòng theo thứ tự tăng dần.

\usacoproblem{Superprime Rib}{sprime}{Superprime Rib}
\ioheading{Nhiệm vụ.}
Một số nguyên tố đặc biệt độ dài \texttt{N} là số mà mọi tiền tố thu được bằng cách cắt dần chữ số bên phải đều là nguyên tố. Cần liệt kê tất cả số như vậy.

\ioheading{Dữ liệu vào.}
Một dòng chứa \texttt{N} (\texttt{1 <= N <= 8}).

\ioheading{Kết quả.}
In các số nguyên tố đặc biệt độ dài \texttt{N} theo thứ tự tăng dần.

\usacoproblem{Checker Challenge}{checker}{Checker Challenge}
\ioheading{Nhiệm vụ.}
Đặt \texttt{N} quân hậu trên bàn cờ \(N\times N\) sao cho không có hai quân cùng hàng, cột hay đường chéo. Mỗi nghiệm được ghi bằng dãy cột của các quân trên từng hàng.

\ioheading{Dữ liệu vào.}
Một số \texttt{N} với \texttt{6 <= N <= 13}.

\ioheading{Kết quả.}
In ba nghiệm đầu theo thứ tự từ điển, rồi in tổng số nghiệm.

\section*{Chương 2}
\subsection*{Section 2.1}
\usacoproblem{The Castle}{castle}{The Castle}
\ioheading{Nhiệm vụ.}
Bản đồ lâu đài là lưới ô, mỗi ô mã hóa các bức tường ở bốn phía. Cần tìm số phòng, diện tích phòng lớn nhất, rồi chọn một bức tường để phá sao cho phòng mới lớn nhất theo đúng quy tắc ưu tiên của đề.

\ioheading{Dữ liệu vào.}
Dòng đầu chứa số cột và số hàng; mỗi ô là một số mã tường.

\ioheading{Kết quả.}
In số phòng, diện tích lớn nhất, diện tích lớn nhất sau khi phá một tường, và vị trí/hướng của tường cần phá.

\usacoproblem{Ordered Fractions}{frac1}{Ordered Fractions}
\ioheading{Nhiệm vụ.}
Liệt kê mọi phân số tối giản trong đoạn từ 0 đến 1 có mẫu số không vượt quá \texttt{N}.

\ioheading{Dữ liệu vào.}
Một số nguyên \texttt{N}.

\ioheading{Kết quả.}
In các phân số theo thứ tự tăng dần.

\usacoproblem{Sorting a Three-Valued Sequence}{sort3}{Sorting a Three-Valued Sequence}
\ioheading{Nhiệm vụ.}
Cho một dãy chỉ gồm 1, 2, 3. Cần sắp tăng dãy bằng số lần hoán đổi hai phần tử bất kỳ ít nhất.

\ioheading{Dữ liệu vào.}
Dòng đầu là độ dài, các dòng sau là từng phần tử.

\ioheading{Kết quả.}
In số phép đổi nhỏ nhất.

\usacoproblem{Healthy Holsteins}{holstein}{Healthy Holsteins}
\ioheading{Nhiệm vụ.}
Bessie cần đạt tối thiểu một lượng vitamin cho mỗi loại. Mỗi thức ăn cung cấp một vector vitamin. Hãy chọn số thức ăn ít nhất, và nếu hòa chọn danh sách chỉ số nhỏ nhất theo thứ tự từ điển.

\ioheading{Dữ liệu vào.}
Dữ liệu cho số loại vitamin, nhu cầu từng loại, số thức ăn và lượng vitamin của mỗi thức ăn.

\ioheading{Kết quả.}
In số thức ăn chọn và chỉ số của chúng.

\usacoproblem{Hamming Codes}{hamming}{Hamming Codes}
\ioheading{Nhiệm vụ.}
Tìm \texttt{N} mã nhị phân dài \texttt{B} bit sao cho khoảng cách Hamming giữa mọi cặp mã ít nhất là \texttt{D}.

\ioheading{Dữ liệu vào.}
Một dòng chứa \texttt{N}, \texttt{B}, \texttt{D}.

\ioheading{Kết quả.}
In các mã theo thứ tự tăng, mỗi dòng tối đa 10 số.

\subsection*{Section 2.2}
\usacoproblem{Preface Numbering}{preface}{Preface Numbering}
\ioheading{Nhiệm vụ.}
Đánh số các trang từ 1 đến \texttt{N} bằng chữ số La Mã. Cần đếm mỗi ký hiệu \texttt{I,V,X,L,C,D,M} xuất hiện bao nhiêu lần.

\ioheading{Dữ liệu vào.}
Một số \texttt{N}.

\ioheading{Kết quả.}
In các ký hiệu có số lần xuất hiện dương cùng số đếm, theo thứ tự giá trị tăng.

\usacoproblem{Subset Sums}{subset}{Subset Sums}
\ioheading{Nhiệm vụ.}
Đếm số cách chia tập \(\{1,2,\ldots,N\}\) thành hai tập con có tổng bằng nhau.

\ioheading{Dữ liệu vào.}
Một số nguyên \texttt{N}.

\ioheading{Kết quả.}
In số cách chia, trong đó hai tập con đổi chỗ được coi là cùng một cách.

\usacoproblem{Runaround Numbers}{runround}{Runaround Numbers}
\ioheading{Nhiệm vụ.}
Một số runaround có các chữ số khác 0 và không lặp; bắt đầu từ chữ số đầu, mỗi bước nhảy theo giá trị chữ số hiện tại modulo độ dài, và phải thăm mọi chữ số đúng một lần rồi quay về đầu.

\ioheading{Dữ liệu vào.}
Một số nguyên \texttt{M}.

\ioheading{Kết quả.}
In số runaround nhỏ nhất lớn hơn \texttt{M}.

\usacoproblem{Party Lamps}{lamps}{Party Lamps}
\ioheading{Nhiệm vụ.}
Có \texttt{N} bóng đèn ban đầu đều bật. Bốn nút sẽ đảo trạng thái các tập bóng: tất cả, vị trí lẻ, vị trí chẵn, và vị trí dạng \(3k+1\). Sau đúng hoặc không quá \texttt{C} lần bấm theo quy tắc của đề, một số bóng bắt buộc bật/tắt.

\ioheading{Dữ liệu vào.}
Dữ liệu gồm \texttt{N}, \texttt{C}, danh sách bóng phải bật và danh sách bóng phải tắt, mỗi danh sách kết thúc bằng \texttt{-1}.

\ioheading{Kết quả.}
In mọi trạng thái có thể theo thứ tự từ điển; nếu không có thì in \texttt{IMPOSSIBLE}.

\subsection*{Section 2.3}
\usacoproblem{Longest Prefix}{prefix}{Longest Prefix}
\ioheading{Nhiệm vụ.}
Cho một tập các primitive và một chuỗi dài. Cần tìm độ dài tiền tố dài nhất của chuỗi có thể ghép từ các primitive.

\ioheading{Dữ liệu vào.}
Các primitive xuất hiện trước dấu chấm, sau đó là chuỗi cần xét.

\ioheading{Kết quả.}
In độ dài tiền tố dài nhất.

\usacoproblem{Cow Pedigrees}{nocows}{Cow Pedigrees}
\ioheading{Nhiệm vụ.}
Đếm số cây nhị phân đầy đủ có đúng \texttt{N} nút và chiều cao đúng \texttt{K}.

\ioheading{Dữ liệu vào.}
Một dòng chứa \texttt{N} và \texttt{K}.

\ioheading{Kết quả.}
In số cây modulo 9901.

\usacoproblem{Zero Sum}{zerosum}{Zero Sum}
\ioheading{Nhiệm vụ.}
Chèn giữa các số 1 đến \texttt{N} một trong ba ký hiệu: khoảng trắng để nối số, \texttt{+}, hoặc \texttt{-}. Cần liệt kê các biểu thức có giá trị bằng 0.

\ioheading{Dữ liệu vào.}
Một số \texttt{N}.

\ioheading{Kết quả.}
In các biểu thức hợp lệ theo thứ tự ASCII của ký hiệu.

\usacoproblem{Money Systems}{money}{Money Systems}
\ioheading{Nhiệm vụ.}
Cho các mệnh giá tiền. Đếm số cách tạo tổng \texttt{N} bằng số lượng không giới hạn các đồng tiền đó.

\ioheading{Dữ liệu vào.}
Dữ liệu chứa số mệnh giá \texttt{V}, tổng cần tạo \texttt{N}, rồi danh sách mệnh giá.

\ioheading{Kết quả.}
In số cách tạo tổng.

\usacoproblem{Controlling Companies}{concom}{Controlling Companies}
\ioheading{Nhiệm vụ.}
Một công ty kiểm soát công ty khác nếu nó trực tiếp hoặc gián tiếp sở hữu hơn 50 phần trăm cổ phần. Cần tìm mọi cặp công ty kiểm soát nhau theo dữ liệu sở hữu ban đầu.

\ioheading{Dữ liệu vào.}
Dòng đầu là số quan hệ; mỗi dòng sau gồm công ty sở hữu, công ty bị sở hữu và phần trăm.

\ioheading{Kết quả.}
In các cặp \(i,j\) với \(i\) kiểm soát \(j\), theo thứ tự tăng.

\subsection*{Section 2.4}
\usacoproblem{The Tamworth Two}{ttwo}{The Tamworth Two}
\ioheading{Nhiệm vụ.}
Farmer John và đàn bò di chuyển trong lưới 10x10 có chướng ngại. Mỗi phút, nếu ô trước mặt trống thì tiến lên, ngược lại quay phải. Cần biết sau bao lâu họ gặp nhau, hoặc không bao giờ gặp.

\ioheading{Dữ liệu vào.}
Mười dòng bản đồ gồm ô trống, chướng ngại, vị trí Farmer John và bò.

\ioheading{Kết quả.}
In số phút tới khi gặp nhau, hoặc 0 nếu chu kỳ không gặp.

\usacoproblem{Overfencing}{maze1}{Overfencing}
\ioheading{Nhiệm vụ.}
Cho mê cung bằng ký tự ASCII có một hoặc hai lối ra. Cần tìm khoảng cách ngắn nhất từ mỗi ô tới một lối ra và lấy giá trị lớn nhất.

\ioheading{Dữ liệu vào.}
Dòng đầu chứa chiều rộng và chiều cao theo số ô; sau đó là bản vẽ mê cung kích thước ký tự.

\ioheading{Kết quả.}
In số bước lớn nhất cần đi để thoát mê cung.

\usacoproblem{Bessie Come Home}{comehome}{Bessie Come Home}
\ioheading{Nhiệm vụ.}
Các đồng cỏ được đặt tên bằng chữ cái. Bessie ở đồng cỏ \texttt{Z}; cần tìm con bò ở đồng cỏ chữ hoa khác \texttt{Z} có đường đi ngắn nhất về \texttt{Z}.

\ioheading{Dữ liệu vào.}
Dòng đầu là số đường; mỗi dòng sau gồm hai tên đồng cỏ và độ dài đường hai chiều.

\ioheading{Kết quả.}
In tên đồng cỏ gần nhất và khoảng cách.

\usacoproblem{Cow Tours}{cowtour}{Cow Tours}
\ioheading{Nhiệm vụ.}
Một số đồng cỏ tạo thành nhiều thành phần liên thông. Có thể thêm đúng một đường nối giữa hai đồng cỏ thuộc hai thành phần khác nhau. Cần đường kính nhỏ nhất có thể của toàn bộ hệ thống sau khi nối.

\ioheading{Dữ liệu vào.}
Dữ liệu gồm số đồng cỏ, tọa độ từng đồng cỏ và ma trận kề.

\ioheading{Kết quả.}
In đường kính nhỏ nhất với sáu chữ số sau dấu thập phân.

\usacoproblem{Fractions to Decimals}{fracdec}{Fractions to Decimals}
\ioheading{Nhiệm vụ.}
Chuyển phân số \(N/D\) thành dạng thập phân; phần tuần hoàn nếu có phải được đặt trong ngoặc đơn.

\ioheading{Dữ liệu vào.}
Một dòng chứa tử số và mẫu số.

\ioheading{Kết quả.}
In biểu diễn thập phân, xuống dòng sao cho mỗi dòng không quá 76 ký tự.

\section*{Chương 3}
\subsection*{Section 3.1}
\usacoproblem{Agri-Net}{agrinet}{Agri-Net}
\ioheading{Nhiệm vụ.}
Cho ma trận chi phí nối mạng giữa các trang trại. Cần tổng chi phí nhỏ nhất để mọi trang trại liên thông.

\ioheading{Dữ liệu vào.}
Dòng đầu là \texttt{N}; sau đó là ma trận \texttt{N x N} chi phí.

\ioheading{Kết quả.}
In trọng số cây khung nhỏ nhất.

\usacoproblem{Score Inflation}{inflate}{Score Inflation}
\ioheading{Nhiệm vụ.}
Một cuộc thi có thời lượng giới hạn. Mỗi loại bài có số điểm và thời gian, có thể làm nhiều bài cùng loại. Cần tối đa hóa điểm.

\ioheading{Dữ liệu vào.}
Dòng đầu chứa tổng thời gian và số loại bài; mỗi dòng sau là điểm và thời gian của một loại.

\ioheading{Kết quả.}
In điểm tối đa.

\usacoproblem{Humble Numbers}{humble}{Humble Numbers}
\ioheading{Nhiệm vụ.}
Cho \texttt{K} số nguyên tố. Một humble number là số có mọi thừa số nguyên tố thuộc tập này. Cần tìm humble number thứ \texttt{N}.

\ioheading{Dữ liệu vào.}
Dữ liệu chứa \texttt{K}, \texttt{N} và danh sách \texttt{K} số nguyên tố.

\ioheading{Kết quả.}
In humble number thứ \texttt{N}, với 1 không được tính là số thứ nhất trong ngữ cảnh kết quả của bài.

\usacoproblem{Shaping Regions}{rect1}{Shaping Regions}
\ioheading{Nhiệm vụ.}
Một hình chữ nhật lớn ban đầu có màu 1. Nhiều hình chữ nhật màu được vẽ chồng lên nhau theo thứ tự. Cần tính diện tích cuối cùng của từng màu còn nhìn thấy.

\ioheading{Dữ liệu vào.}
Dữ liệu gồm kích thước nền, số hình chữ nhật, rồi tọa độ và màu của từng hình.

\ioheading{Kết quả.}
In mỗi màu có diện tích dương cùng diện tích của nó, theo màu tăng dần.

\usacoproblem{Contact}{contact}{Contact}
\ioheading{Nhiệm vụ.}
Trong một chuỗi bit dài, đếm tần suất mọi mẫu nhị phân có độ dài từ \texttt{A} đến \texttt{B}. Cần in \texttt{N} nhóm tần suất cao nhất.

\ioheading{Dữ liệu vào.}
Dòng đầu chứa \texttt{A}, \texttt{B}, \texttt{N}; phần còn lại là chuỗi bit.

\ioheading{Kết quả.}
Với mỗi tần suất được chọn, in tần suất rồi các mẫu có tần suất đó theo độ dài và giá trị tăng dần.

\usacoproblem{Stamps}{stamps}{Stamps}
\ioheading{Nhiệm vụ.}
Có thể dán tối đa \texttt{K} con tem lên một phong bì, các mệnh giá tem cho trước và dùng không giới hạn. Hỏi mọi giá trị từ 1 liên tiếp đến giá trị lớn nhất nào đều có thể tạo.

\ioheading{Dữ liệu vào.}
Dữ liệu chứa \texttt{K}, số mệnh giá và danh sách mệnh giá.

\ioheading{Kết quả.}
In giá trị lớn nhất trong đoạn liên tiếp bắt đầu từ 1 có thể tạo.

\subsection*{Section 3.2}
\usacoproblem{Factorials}{fact4}{Factorials}
\ioheading{Nhiệm vụ.}
Tìm chữ số khác 0 cuối cùng của \(N!\).

\ioheading{Dữ liệu vào.}
Một số nguyên \texttt{N}.

\ioheading{Kết quả.}
In chữ số khác 0 cuối cùng.

\usacoproblem{Stringsobits}{kimbits}{Stringsobits}
\ioheading{Nhiệm vụ.}
Xét các chuỗi bit độ dài \texttt{N} có không quá \texttt{L} bit 1, sắp theo thứ tự từ điển. Cần tìm chuỗi thứ \texttt{I}.

\ioheading{Dữ liệu vào.}
Một dòng chứa \texttt{N}, \texttt{L}, \texttt{I}.

\ioheading{Kết quả.}
In chuỗi bit cần tìm.

\usacoproblem{Spinning Wheels}{spin}{Spinning Wheels}
\ioheading{Nhiệm vụ.}
Năm bánh xe quay, mỗi bánh có tốc độ và một số khe hình nêm. Cần tìm thời điểm đầu tiên trong 360 giây đầu mà có một góc đi xuyên qua khe của cả năm bánh.

\ioheading{Dữ liệu vào.}
Mỗi dòng mô tả tốc độ của một bánh, số khe, rồi vị trí bắt đầu và độ rộng của từng khe.

\ioheading{Kết quả.}
In thời điểm nhỏ nhất, hoặc \texttt{none} nếu không có.

\usacoproblem{Feed Ratios}{ratios}{Feed Ratios}
\ioheading{Nhiệm vụ.}
Ba loại thức ăn có tỉ lệ thành phần khác nhau. Cần trộn các lượng nguyên không âm để đạt một bội nguyên dương của tỉ lệ mục tiêu, với tổng lượng nhỏ nhất.

\ioheading{Dữ liệu vào.}
Dòng đầu là tỉ lệ mục tiêu; ba dòng sau là thành phần của từng loại thức ăn.

\ioheading{Kết quả.}
In ba lượng thức ăn và hệ số bội; nếu không thể thì in \texttt{NONE}.

\usacoproblem{Magic Squares}{msquare}{Magic Squares}
\ioheading{Nhiệm vụ.}
Trạng thái gồm 8 số trên một bảng 2x4. Ba phép biến đổi A, B, C cho trước thay đổi thứ tự các số. Cần biến trạng thái chuẩn thành trạng thái đích bằng dãy phép ngắn nhất, hòa thì chọn thứ tự từ điển nhỏ nhất.

\ioheading{Dữ liệu vào.}
Một dòng chứa hoán vị đích của các số 1 đến 8.

\ioheading{Kết quả.}
In độ dài dãy phép và dãy ký tự phép biến đổi.

\usacoproblem{Sweet Butter}{butter}{Sweet Butter}
\ioheading{Nhiệm vụ.}
Các con bò đứng ở một số đồng cỏ trong đồ thị có trọng số. Chọn đồng cỏ đặt bơ sao cho tổng khoảng cách các bò đi tới đó là nhỏ nhất.

\ioheading{Dữ liệu vào.}
Dữ liệu gồm số bò, số đồng cỏ, số đường; sau đó là vị trí bò và danh sách đường hai chiều có độ dài.

\ioheading{Kết quả.}
In tổng khoảng cách nhỏ nhất.

\subsection*{Section 3.3}
\usacoproblem{Riding the Fences}{fence}{Riding the Fences}
\ioheading{Nhiệm vụ.}
Cho một đa đồ thị vô hướng biểu diễn các đoạn hàng rào. Cần tìm hành trình đi qua mỗi đoạn đúng một lần, bắt đầu ở đỉnh nhỏ nhất hợp lệ và có thứ tự từ điển nhỏ nhất.

\ioheading{Dữ liệu vào.}
Dòng đầu là số cạnh, sau đó là các cặp đỉnh.

\ioheading{Kết quả.}
In các đỉnh của đường đi Euler, mỗi đỉnh một dòng.

\usacoproblem{Shopping Offers}{shopping}{Shopping Offers}
\ioheading{Nhiệm vụ.}
Có các gói khuyến mãi, mỗi gói gồm một số mặt hàng và giá. Với danh sách cần mua tối đa năm loại hàng, tính chi phí nhỏ nhất.

\ioheading{Dữ liệu vào.}
Dữ liệu mô tả các gói khuyến mãi, rồi các mặt hàng cần mua, số lượng và giá thường.

\ioheading{Kết quả.}
In chi phí nhỏ nhất để mua đúng số lượng cần.

\usacoproblem{Camelot}{camelot}{Camelot}
\ioheading{Nhiệm vụ.}
Trên bàn cờ có một vua và nhiều mã. Tất cả quân phải tập hợp tại một ô với tổng số bước nhỏ nhất; một mã có thể đón vua và cùng đi tiếp, bước đi của cặp được tính như bước của mã.

\ioheading{Dữ liệu vào.}
Dòng đầu chứa số hàng và cột. Các cặp chữ-cố chỉ vị trí vua trước, sau đó là các mã.

\ioheading{Kết quả.}
In tổng số bước nhỏ nhất.

\usacoproblem{Home on the Range}{range}{Home on the Range}
\ioheading{Nhiệm vụ.}
Một đồng cỏ vuông được chia thành ô tốt/xấu. Cần đếm số hình vuông toàn ô tốt theo từng kích thước từ 2 trở lên.

\ioheading{Dữ liệu vào.}
Dòng đầu là \texttt{N}; tiếp theo là \texttt{N} dòng bit.

\ioheading{Kết quả.}
Với mỗi kích thước có ít nhất một hình vuông, in kích thước và số lượng.

\usacoproblem{A Game}{game1}{A Game}
\ioheading{Nhiệm vụ.}
Hai người chơi lần lượt lấy một số ở một trong hai đầu dãy. Cả hai chơi tối ưu. Cần tính điểm cuối cùng của hai người.

\ioheading{Dữ liệu vào.}
Dòng đầu là \texttt{N}; sau đó là \texttt{N} số của dãy.

\ioheading{Kết quả.}
In điểm của người đi trước và người đi sau.

\subsection*{Section 3.4}
\usacoproblem{Closed Fences}{fence4}{Closed Fences}
\ioheading{Nhiệm vụ.}
Một hàng rào đa giác khép kín được quan sát từ một điểm. Cần xác định những đoạn hàng rào có ít nhất một phần nhìn thấy từ điểm quan sát.

\ioheading{Dữ liệu vào.}
Dữ liệu gồm điểm quan sát, số đỉnh và tọa độ các đỉnh của hàng rào.

\ioheading{Kết quả.}
In số đoạn nhìn thấy và danh sách các đoạn theo thứ tự quy định trong đề.

\usacoproblem{American Heritage}{heritage}{American Heritage}
\ioheading{Nhiệm vụ.}
Cho thứ tự duyệt trung tự và tiền tự của một cây nhị phân với nhãn chữ cái khác nhau. Cần khôi phục và in thứ tự hậu tự.

\ioheading{Dữ liệu vào.}
Hai dòng lần lượt là chuỗi trung tự và tiền tự.

\ioheading{Kết quả.}
In chuỗi hậu tự.

\usacoproblem{Electric Fences}{fence3}{Electric Fences}
\ioheading{Nhiệm vụ.}
Cho các đoạn hàng rào song song trục tọa độ. Cần chọn vị trí đặt nguồn điện sao cho tổng khoảng cách ngắn nhất từ nguồn tới từng đoạn là nhỏ nhất.

\ioheading{Dữ liệu vào.}
Dòng đầu là số đoạn; mỗi dòng sau chứa tọa độ hai đầu đoạn.

\ioheading{Kết quả.}
In tọa độ tốt nhất và tổng độ dài dây, làm tròn theo yêu cầu.

\usacoproblem{Raucous Rockers}{rockers}{Raucous Rockers}
\ioheading{Nhiệm vụ.}
Có \texttt{N} bài hát theo thứ tự, mỗi đĩa chứa tối đa \texttt{T} phút và có \texttt{M} đĩa. Chọn nhiều bài nhất có thể, giữ nguyên thứ tự ban đầu.

\ioheading{Dữ liệu vào.}
Dữ liệu chứa \texttt{N}, \texttt{T}, \texttt{M} và thời lượng từng bài.

\ioheading{Kết quả.}
In số bài tối đa có thể ghi.

\section*{Chương 4}
\subsection*{Section 4.1}
\usacoproblem{Beef McNuggets}{nuggets}{Beef McNuggets}
\ioheading{Nhiệm vụ.}
Cho các kích cỡ gói McNuggets. Tìm số nguyên lớn nhất không thể biểu diễn bằng tổng không âm các kích cỡ gói; nếu có vô hạn số không biểu diễn được thì in 0.

\ioheading{Dữ liệu vào.}
Dòng đầu là số kích cỡ, sau đó là từng kích cỡ.

\ioheading{Kết quả.}
In số lớn nhất không biểu diễn được, hoặc 0.

\usacoproblem{Fence Rails}{fence8}{Fence Rails}
\ioheading{Nhiệm vụ.}
Có các tấm ván dài và các thanh ray cần cắt. Hãy cắt được nhiều thanh ray nhất có thể, mỗi thanh từ một tấm ván, phần thừa có thể bỏ.

\ioheading{Dữ liệu vào.}
Dữ liệu cho số tấm ván và độ dài của chúng, rồi số ray và độ dài từng ray.

\ioheading{Kết quả.}
In số ray nhiều nhất có thể cắt.

\usacoproblem{Fence Loops}{fence6}{Fence Loops}
\ioheading{Nhiệm vụ.}
Mỗi đoạn hàng rào biết độ dài và các đoạn nối vào hai đầu của nó. Cần tìm chu trình hàng rào có tổng độ dài nhỏ nhất.

\ioheading{Dữ liệu vào.}
Dữ liệu mô tả từng đoạn: mã đoạn, độ dài, số đoạn kề ở hai đầu và danh sách các đoạn kề.

\ioheading{Kết quả.}
In độ dài chu trình nhỏ nhất.

\usacoproblem{Cryptcowgraphy}{cryptcow}{Cryptcowgraphy}
\ioheading{Nhiệm vụ.}
Một thông điệp đã mã hóa bằng cách lặp thao tác chèn ba chữ cái \texttt{C}, \texttt{O}, \texttt{W} quanh ba đoạn của chuỗi. Cần kiểm tra có thể giải mã về câu gốc chuẩn hay không và số bước.

\ioheading{Dữ liệu vào.}
Một dòng là thông điệp cần kiểm tra.

\ioheading{Kết quả.}
In \texttt{1} và số bước nếu hợp lệ, ngược lại in \texttt{0 0}.

\subsection*{Section 4.2}
\usacoproblem{Drainage Ditches}{ditch}{Drainage Ditches}
\ioheading{Nhiệm vụ.}
Cho mạng mương thoát nước có hướng với dung lượng. Cần lượng nước tối đa chảy từ nút 1 tới nút \texttt{M}.

\ioheading{Dữ liệu vào.}
Dòng đầu chứa số cạnh và số nút; mỗi dòng sau là cạnh có hướng và dung lượng.

\ioheading{Kết quả.}
In giá trị luồng cực đại.

\usacoproblem{The Perfect Stall}{stall4}{The Perfect Stall}
\ioheading{Nhiệm vụ.}
Mỗi con bò chấp nhận một số chuồng. Cần xếp nhiều bò nhất vào các chuồng khác nhau mà chúng chấp nhận.

\ioheading{Dữ liệu vào.}
Dữ liệu gồm số bò, số chuồng, rồi danh sách chuồng chấp nhận của từng bò.

\ioheading{Kết quả.}
In số bò tối đa được xếp.

\usacoproblem{Job Processing}{job}{Job Processing}
\ioheading{Nhiệm vụ.}
Có \texttt{N} công việc, mỗi công việc cần qua nhóm máy A rồi nhóm máy B. Các máy trong cùng nhóm có thời gian xử lý khác nhau. Cần thời điểm xong toàn bộ giai đoạn A và thời điểm hoàn thành tất cả.

\ioheading{Dữ liệu vào.}
Dữ liệu chứa số công việc, số máy A, số máy B, rồi thời gian của từng máy.

\ioheading{Kết quả.}
In hai thời điểm theo yêu cầu.

\usacoproblem{Cowcycles}{cowcycle}{Cowcycles}
\ioheading{Nhiệm vụ.}
Chọn các đĩa trước và sau cho xe đạp sao cho tỉ số truyền nằm trong miền cho phép và phương sai của các hiệu giữa tỉ số liên tiếp là nhỏ nhất.

\ioheading{Dữ liệu vào.}
Dữ liệu cho số đĩa cần chọn và khoảng giá trị cho đĩa trước/sau.

\ioheading{Kết quả.}
In bộ đĩa trước và bộ đĩa sau tối ưu theo thứ tự tăng.

\subsection*{Section 4.3}
\usacoproblem{Buy Low, Buy Lower}{buylow}{Buy Low, Buy Lower}
\ioheading{Nhiệm vụ.}
Cho dãy giá cổ phiếu. Cần độ dài dãy con giảm dài nhất và số dãy con giảm dài nhất khác nhau, không đếm trùng các dãy có cùng giá trị.

\ioheading{Dữ liệu vào.}
Dòng đầu là \texttt{N}; sau đó là dãy giá.

\ioheading{Kết quả.}
In độ dài và số lượng dãy con.

\usacoproblem{The Primes}{prime3}{The Primes}
\ioheading{Nhiệm vụ.}
Điền các chữ số vào bảng 5x5 sao cho mỗi hàng, cột và hai đường chéo chính tạo thành số nguyên tố năm chữ số có cùng tổng chữ số; chữ số góc trên trái được cho trước.

\ioheading{Dữ liệu vào.}
Một dòng chứa tổng chữ số và chữ số đầu tiên.

\ioheading{Kết quả.}
In tất cả bảng hợp lệ theo thứ tự từ điển, mỗi bảng năm dòng; nếu không có thì in \texttt{NONE}.

\usacoproblem{Street Race}{race3}{Street Race}
\ioheading{Nhiệm vụ.}
Trong một đồ thị có hướng tốt từ đỉnh đầu đến đỉnh cuối, tìm các đỉnh không thể tránh khỏi trên mọi đường đi và các đỉnh có thể tách đường đua thành hai phần độc lập.

\ioheading{Dữ liệu vào.}
Mỗi dòng liệt kê các cung đi ra của một đỉnh, kết thúc bằng \texttt{-2}; toàn bộ dữ liệu kết thúc bằng \texttt{-1}.

\ioheading{Kết quả.}
In hai dòng: danh sách đỉnh không thể tránh khỏi và danh sách đỉnh tách, mỗi dòng bắt đầu bằng số lượng.

\usacoproblem{Letter Game}{lgame}{Letter Game}
\ioheading{Nhiệm vụ.}
Từ các chữ cái thu thập được, chọn một từ hoặc một cặp từ trong từ điển để đạt tổng điểm chữ cái lớn nhất, không dùng quá số chữ cái có sẵn.

\ioheading{Dữ liệu vào.}
Tệp \texttt{lgame.in} chứa các chữ cái; tệp \texttt{lgame.dict} chứa từ điển kết thúc bằng dấu chấm.

\ioheading{Kết quả.}
In điểm lớn nhất rồi mọi từ hoặc cặp từ đạt điểm đó, theo thứ tự yêu cầu.

\subsection*{Section 4.4}
\usacoproblem{Shuttle Puzzle}{shuttle}{Shuttle Puzzle}
\ioheading{Nhiệm vụ.}
Có \texttt{N} quân trắng, \texttt{N} quân đen và một ô trống trên hàng dài \(2N+1\). Quân có thể đi sang ô trống kề bên hoặc nhảy qua đúng một quân khác màu. Cần đổi chỗ hai màu bằng dãy bước ít nhất, hòa chọn từ điển nhỏ nhất.

\ioheading{Dữ liệu vào.}
Một số \texttt{N}.

\ioheading{Kết quả.}
In các vị trí được di chuyển, mỗi dòng tối đa 20 số.

\usacoproblem{Pollutant Control}{milk6}{Pollutant Control}
\ioheading{Nhiệm vụ.}
Trong mạng vận chuyển sữa có hướng, cần chặn mọi đường từ nhà máy đến cửa hàng bằng cách dừng một số xe tải. Tối thiểu hóa tổng chi phí, rồi số xe, rồi danh sách chỉ số theo thứ tự từ điển.

\ioheading{Dữ liệu vào.}
Dòng đầu là số tuyến và số kho; mỗi dòng sau là tuyến có hướng cùng chi phí dừng.

\ioheading{Kết quả.}
In chi phí nhỏ nhất và số tuyến dừng, rồi chỉ số các tuyến cần dừng.

\usacoproblem{Frame Up}{frameup}{Frame Up}
\ioheading{Nhiệm vụ.}
Một bức ảnh gồm nhiều khung chữ nhật bằng chữ cái bị chồng lên nhau. Cần liệt kê mọi thứ tự có thể dùng để đặt các khung từ dưới lên tạo ra ảnh cuối.

\ioheading{Dữ liệu vào.}
Dữ liệu là kích thước ảnh và các dòng ký tự, dấu chấm là nền.

\ioheading{Kết quả.}
In mọi thứ tự khung hợp lệ theo thứ tự từ điển.

\section*{Chương 5}
\subsection*{Section 5.1}
\usacoproblem{Fencing the Cows}{fc}{Fencing the Cows}
\ioheading{Nhiệm vụ.}
Cho tọa độ các con bò. Cần chiều dài hàng rào ngắn nhất bao quanh tất cả bò, tức chu vi bao lồi.

\ioheading{Dữ liệu vào.}
Dòng đầu là số điểm, sau đó là tọa độ từng điểm.

\ioheading{Kết quả.}
In chu vi với hai chữ số sau dấu thập phân.

\usacoproblem{Starry Night}{starry}{Starry Night}
\ioheading{Nhiệm vụ.}
Bầu trời là lưới bit. Các cụm sao liên thông tám hướng được coi là cùng loại nếu có thể trùng nhau qua tịnh tiến, quay hoặc phản chiếu. Cần gán cùng chữ cái cho các cụm cùng loại.

\ioheading{Dữ liệu vào.}
Dữ liệu gồm chiều rộng, chiều cao và bản đồ bit.

\ioheading{Kết quả.}
In bản đồ mới với chữ cái cho sao và \texttt{0} cho nền.

\usacoproblem{Musical Themes}{theme}{Musical Themes}
\ioheading{Nhiệm vụ.}
Một theme là hai đoạn con không chồng nhau có cùng chuỗi hiệu cao độ liên tiếp và độ dài nốt ít nhất 5. Cần tìm độ dài theme dài nhất.

\ioheading{Dữ liệu vào.}
Dòng đầu là số nốt, sau đó là cao độ các nốt.

\ioheading{Kết quả.}
In độ dài lớn nhất, hoặc 0 nếu không có theme đủ dài.

\subsection*{Section 5.2}
\usacoproblem{Snail Trails}{snail}{Snail Trails}
\ioheading{Nhiệm vụ.}
Một con ốc đi trong lưới vuông có chướng ngại. Nó trượt thẳng đến khi gặp mép hoặc chướng ngại, rồi có thể rẽ trái/phải. Không được đi qua ô đã thăm. Cần số ô tối đa có thể thăm.

\ioheading{Dữ liệu vào.}
Dữ liệu gồm kích thước lưới, số chướng ngại và tọa độ các chướng ngại.

\ioheading{Kết quả.}
In số ô lớn nhất có thể đi qua.

\usacoproblem{Electric Fences}{fence3}{Electric Fences}
\ioheading{Nhiệm vụ.}
Bản PDF lặp lại bài \textit{Electric Fences}: chọn vị trí nguồn điện tối ưu để tổng khoảng cách tới các đoạn hàng rào song song trục là nhỏ nhất.

\ioheading{Dữ liệu vào.}
Dòng đầu là số đoạn; mỗi dòng sau chứa hai đầu mút của đoạn.

\ioheading{Kết quả.}
In tọa độ và tổng khoảng cách tối thiểu.

\usacoproblem{Wisconsin Squares}{wissqu}{Wisconsin Squares}
\ioheading{Nhiệm vụ.}
Trên bảng 4x4 có các chữ \texttt{A} đến \texttt{E}. Cần thay toàn bộ 16 ô bằng một dãy chữ cái theo số lượng cho trước, sao cho mỗi chữ đặt vào không kề tám hướng với cùng chữ đó trong trạng thái hiện tại.

\ioheading{Dữ liệu vào.}
Bốn dòng mô tả bảng ban đầu.

\ioheading{Kết quả.}
In một dãy đặt hợp lệ theo yêu cầu và tổng số dãy hợp lệ.

\subsection*{Section 5.3}
\usacoproblem{Milk Measuring}{milk4}{Milk Measuring}
\ioheading{Nhiệm vụ.}
Cho các kích cỡ ca đong và một lượng sữa cần đo. Cần chọn ít ca nhất sao cho có thể đo đúng lượng đó bằng cách dùng không giới hạn các ca đã chọn; hòa thì chọn danh sách kích cỡ nhỏ nhất theo thứ tự từ điển.

\ioheading{Dữ liệu vào.}
Dữ liệu gồm lượng cần đo, số ca và kích cỡ từng ca.

\ioheading{Kết quả.}
In số ca chọn và kích cỡ của chúng.

\usacoproblem{Window Area}{window}{Window Area}
\ioheading{Nhiệm vụ.}
Mô phỏng các cửa sổ chữ nhật chồng lớp. Lệnh có thể tạo cửa sổ, đưa lên trên, đưa xuống dưới, xóa, hoặc hỏi phần trăm diện tích nhìn thấy của một cửa sổ.

\ioheading{Dữ liệu vào.}
Mỗi dòng là một lệnh dạng \texttt{w}, \texttt{t}, \texttt{b}, \texttt{d}, hoặc \texttt{s}; dữ liệu kết thúc ở EOF.

\ioheading{Kết quả.}
Với mỗi lệnh \texttt{s}, in phần trăm nhìn thấy làm tròn ba chữ số sau dấu thập phân.

\usacoproblem{Network of Schools}{schlnet}{Network of Schools}
\ioheading{Nhiệm vụ.}
Các trường tạo thành đồ thị có hướng phân phối phần mềm. Cần biết tối thiểu bao nhiêu trường phải nhận phần mềm ban đầu để mọi trường nhận được, và tối thiểu bao nhiêu liên kết cần thêm để từ bất kỳ trường nào phần mềm cũng tới được mọi trường.

\ioheading{Dữ liệu vào.}
Dòng đầu là số trường; mỗi dòng sau liệt kê các trường nhận từ trường đó, kết thúc bằng 0.

\ioheading{Kết quả.}
In hai số trên hai dòng theo hai câu hỏi.

\usacoproblem{Big Barn}{bigbrn}{Big Barn}
\ioheading{Nhiệm vụ.}
Trong một lưới \(N\times N\) có một số ô bị cây chiếm. Cần hình vuông lớn nhất không chứa ô bị chiếm.

\ioheading{Dữ liệu vào.}
Dòng đầu chứa \texttt{N} và số cây; mỗi dòng sau là tọa độ một cây.

\ioheading{Kết quả.}
In cạnh của hình vuông lớn nhất.

\subsection*{Section 5.4}
\usacoproblem{All Latin Squares}{latin}{All Latin Squares}
\ioheading{Nhiệm vụ.}
Đếm số Latin square bậc \texttt{N}, trong đó mỗi hàng và mỗi cột là hoán vị của các ký hiệu, với chuẩn hóa theo hàng đầu để tránh đếm đối xứng thừa.

\ioheading{Dữ liệu vào.}
Một số \texttt{N}.

\ioheading{Kết quả.}
In số Latin square theo quy ước của đề.

\usacoproblem{Canada Tour}{tour}{Canada Tour}
\ioheading{Nhiệm vụ.}
Có các thành phố theo thứ tự đông-tây và các chuyến bay hai chiều chỉ xét như cạnh giữa thành phố. Cần chuyến du lịch đi từ thành phố đầu đến cuối rồi về lại đầu, thăm nhiều thành phố khác nhau nhất có thể theo ràng buộc thứ tự.

\ioheading{Dữ liệu vào.}
Dữ liệu gồm số thành phố, số chuyến bay, tên thành phố theo thứ tự và danh sách cạnh.

\ioheading{Kết quả.}
In số thành phố nhiều nhất có thể thăm.

\usacoproblem{Character Recognition}{char}{Character Recognition}
\ioheading{Nhiệm vụ.}
Nhận dạng văn bản từ ảnh bitmap. Mỗi ký tự chuẩn cao 20 dòng; ảnh đầu vào có thể bị mất hoặc thừa một dòng trong mỗi ký tự. Cần khôi phục chuỗi có sai khác nhỏ nhất.

\ioheading{Dữ liệu vào.}
Dữ liệu dùng tệp phông chuẩn và tệp ảnh đầu vào gồm số dòng rồi các dòng bit.

\ioheading{Kết quả.}
In chuỗi ký tự được nhận dạng.

\usacoproblem{Betsy's Tour}{betsy}{Betsy's Tour}
\ioheading{Nhiệm vụ.}
Trong lưới \(N\times N\), Bessie bắt đầu ở góc trên trái và phải đi tới góc dưới trái, thăm mỗi ô đúng một lần. Cần đếm số đường đi.

\ioheading{Dữ liệu vào.}
Một số \texttt{N}.

\ioheading{Kết quả.}
In số đường đi hợp lệ.

\usacoproblem{Telecowmunication}{telecow}{Telecowmunication}
\ioheading{Nhiệm vụ.}
Trong mạng máy tính vô hướng, cần làm hỏng ít máy tính nhất để hai máy \texttt{c1} và \texttt{c2} không còn liên lạc được. Hai máy đầu cuối không được hỏng; nếu có nhiều tập tối ưu, chọn danh sách từ điển nhỏ nhất.

\ioheading{Dữ liệu vào.}
Dòng đầu chứa \texttt{N}, \texttt{M}, \texttt{c1}, \texttt{c2}; sau đó là \texttt{M} cạnh vô hướng.

\ioheading{Kết quả.}
In số máy cần hỏng và danh sách máy đó.

\subsection*{Section 5.5}
\usacoproblem{Picture}{picture}{Picture}
\ioheading{Nhiệm vụ.}
Cho nhiều hình chữ nhật song song trục có thể chồng lên nhau. Cần chu vi đường bao của hợp các hình chữ nhật.

\ioheading{Dữ liệu vào.}
Dòng đầu là số hình chữ nhật; mỗi dòng sau chứa tọa độ góc dưới trái và góc trên phải.

\ioheading{Kết quả.}
In chu vi của hợp hình.

\usacoproblem{Hidden Password}{hidden}{Hidden Password}
\ioheading{Nhiệm vụ.}
Cho chuỗi vòng tròn độ dài \texttt{L}. Xét mọi phép xoay chuỗi và sắp xếp từ điển; mật khẩu là chỉ số bắt đầu, tính từ 0, của xoay nhỏ nhất.

\ioheading{Dữ liệu vào.}
Dòng đầu là \texttt{L}; dòng sau là chuỗi chữ thường, có thể được ngắt dòng trong dữ liệu.

\ioheading{Kết quả.}
In chỉ số bắt đầu của xoay nhỏ nhất.

\usacoproblem{Two Five}{twofive}{Two Five}
\ioheading{Nhiệm vụ.}
Một từ hợp lệ dùng mỗi chữ \texttt{A} đến \texttt{Y} đúng một lần và khi đặt vào bảng 5x5 thì mỗi hàng, cột đều tăng. Các từ được mã hóa bằng thứ hạng của dãy vị trí hàng-cột của các chữ. Cần đổi qua lại giữa thứ hạng và từ.

\ioheading{Dữ liệu vào.}
Dòng đầu là \texttt{N} để đổi số sang từ hoặc \texttt{W} để đổi từ sang số; dòng sau là số thứ hạng hoặc từ hợp lệ.

\ioheading{Kết quả.}
In từ hoặc số thứ hạng tương ứng.

\section*{Chương 6}
\subsection*{Section 6.1}
\usacoproblem{Postal Vans}{vans}{Postal Vans}
\ioheading{Nhiệm vụ.}
Khu phố có 4 đường đông-tây và \texttt{N} đường bắc-nam. Xe bưu điện xuất phát từ bưu cục ở góc tây bắc và phải đi qua mỗi giao lộ đúng một lần. Cần đếm số tuyến đường.

\ioheading{Dữ liệu vào.}
Một số \texttt{N}.

\ioheading{Kết quả.}
In số tuyến đường.

\usacoproblem{A Rectangular Barn}{rectbarn}{A Rectangular Barn}
\ioheading{Nhiệm vụ.}
Một cánh đồng hình chữ nhật có một số ô 1x1 bị hỏng. Cần diện tích hình chữ nhật trục song song lớn nhất không chứa ô hỏng để xây chuồng.

\ioheading{Dữ liệu vào.}
Dòng đầu chứa số hàng, số cột và số ô hỏng; các dòng sau là tọa độ ô hỏng.

\ioheading{Kết quả.}
In diện tích lớn nhất.

\usacoproblem{Cow XOR}{cowxor}{Cow XOR}
\ioheading{Nhiệm vụ.}
Cho dãy số của các con bò theo thứ tự xã hội. Cần đoạn con liên tiếp có giá trị XOR lớn nhất; nếu hòa chọn đoạn có chỉ số kết thúc nhỏ nhất, rồi đoạn ngắn nhất.

\ioheading{Dữ liệu vào.}
Dòng đầu là \texttt{N}; sau đó là \texttt{N} số trong khoảng từ 0 đến \(2^{21}-1\).

\ioheading{Kết quả.}
In giá trị XOR lớn nhất, chỉ số bắt đầu và chỉ số kết thúc của đoạn.

\clearpage
\section*{Phụ lục: dữ liệu mẫu}
Các khối sau giữ nguyên dữ liệu mẫu và tên tệp mẫu đọc được từ PDF nguồn. Những dòng chú thích tiếng Trung nằm xen giữa ví dụ trong bản trích văn bản đã được bỏ khỏi bản dựng; dữ liệu vào/ra vẫn được giữ dưới dạng nguyên văn.

\paragraph{SAMPLE INPUT (file ride.in)}
\begin{verbatim}
COMETQ
HVNGAT
\end{verbatim}

\paragraph{SAMPLE OUTPUT (file ride.out)}
\begin{verbatim}
GO
\end{verbatim}

\paragraph{SAMPLE INPUT (file gift1.in)}
\begin{verbatim}
5
dave
laura
owen
vick
amr
dave
200 3
laura
owen
vick
owen
500 1
dave
amr
150 2
vick
owen
laura
0 2
amr
vick
vick
0 0
\end{verbatim}

\paragraph{SAMPLE OUTPUT (file gift1.out)}
\begin{verbatim}
dave 302
laura 66
owen -359
vick 141
amr -150
\end{verbatim}

\paragraph{SAMPLE INPUT (file friday.in)}
\begin{verbatim}
20
\end{verbatim}

\paragraph{SAMPLE OUTPUT (file friday.out)}
\begin{verbatim}
36 33 34 33 35 35 34
\end{verbatim}

\paragraph{SAMPLE INPUT (file beads.in)}
\begin{verbatim}
29
wwwbbrwrbrbrrbrbrwrwwrbwrwrrb
\end{verbatim}

\paragraph{SAMPLE OUTPUT (file beads.out)}
\begin{verbatim}
11
\end{verbatim}

\paragraph{SAMPLE INPUT (file milk2.in)}
\begin{verbatim}
3
300 1000
700 1200
1500 2100
\end{verbatim}

\paragraph{SAMPLE OUTPUT (file milk2.out)}
\begin{verbatim}
900 300
\end{verbatim}

\paragraph{SAMPLE INPUT (file transform.in)}
\begin{verbatim}
3
@-@
--@@@-@
@---@
\end{verbatim}

\paragraph{SAMPLE OUTPUT (file transform.out)}
\begin{verbatim}
1
\end{verbatim}

\paragraph{SAMPLE INPUT (file namenum.in)}
\begin{verbatim}
4734
\end{verbatim}

\paragraph{SAMPLE OUTPUT (file namenum.out)}
\begin{verbatim}
GREG
\end{verbatim}

\paragraph{SAMPLE INPUT (file palsquare.in)}
\begin{verbatim}
10
\end{verbatim}

\paragraph{SAMPLE OUTPUT (file palsquare.out)}
\begin{verbatim}
1 1
2 4
3 9
11 121
22 484
26 676
101 10201
111 12321
121 14641
202 40804
212 44944
264 69696
\end{verbatim}

\paragraph{SAMPLE INPUT (file dualpal.in)}
\begin{verbatim}
3 25
\end{verbatim}

\paragraph{SAMPLE OUTPUT (file dualpal.out)}
\begin{verbatim}
26
\end{verbatim}

\paragraph{SAMPLE INPUT (file milk.in)}
\begin{verbatim}
100 5
5 20
9 40
3 10
8 80
6 30
\end{verbatim}

\paragraph{SAMPLE OUTPUT (file milk.out)}
\begin{verbatim}
630
\end{verbatim}

\paragraph{SAMPLE INPUT (file barn1.in)}
\begin{verbatim}
4 50 18
\end{verbatim}

\paragraph{SAMPLE OUTPUT (file barn1.out)}
\begin{verbatim}
25
\end{verbatim}

\paragraph{SAMPLE INPUT (file calfflac.in)}
\begin{verbatim}
Confucius say: Madam, I'm Adam.
\end{verbatim}

\paragraph{SAMPLE OUTPUT (file calfflac.out)}
\begin{verbatim}
11
Madam, I'm Adam
\end{verbatim}

\paragraph{SAMPLE INPUT (file crypt1.in)}
\begin{verbatim}
5
2 3 4 6 8
\end{verbatim}

\paragraph{SAMPLE OUTPUT (file crypt1.out)}
\begin{verbatim}
1
\end{verbatim}

\paragraph{SAMPLE INPUT (file clocks.in)}
\begin{verbatim}
9 9 12
6 6 6
6 3 6
\end{verbatim}

\paragraph{SAMPLE OUTPUT (file clocks.out)}
\begin{verbatim}
4 5 8 9
\end{verbatim}

\paragraph{SAMPLE INPUT (file packrec.in)}
\begin{verbatim}
1 2
2 3
3 4
4 5
\end{verbatim}

\paragraph{SAMPLE OUTPUT (file packrec.out)}
\begin{verbatim}
40
4 10
5 8
\end{verbatim}

\paragraph{SAMPLE INPUT (file ariprog.in)}
\begin{verbatim}
5
\end{verbatim}

\paragraph{SAMPLE OUTPUT (file ariprog.out)}
\begin{verbatim}
1 4
37 4
2 8
29 8
1 12
5 12
13 12
17 12
5 20
2 24
\end{verbatim}

\paragraph{SAMPLE INPUT (file milk3.in)}
\begin{verbatim}
8 9 10
\end{verbatim}

\paragraph{SAMPLE OUTPUT (file milk3.out)}
\begin{verbatim}
1 2 8 9 10
\end{verbatim}

\paragraph{SAMPLE INPUT (file milk3.in)}
\begin{verbatim}
2 5 10
\end{verbatim}

\paragraph{SAMPLE OUTPUT (file milk3.out)}
\begin{verbatim}
5 6 7 8 9 10
\end{verbatim}

\paragraph{SAMPLE INPUT (file numtri.in)}
\begin{verbatim}
5
7
3 8
8 1 0
2 7 4 4
4 5 2 6 5
\end{verbatim}

\paragraph{SAMPLE OUTPUT (file numtri.out)}
\begin{verbatim}
30
\end{verbatim}

\paragraph{SAMPLE INPUT (file pprime.in)}
\begin{verbatim}
5 500
\end{verbatim}

\paragraph{SAMPLE OUTPUT (file pprime.out)}
\begin{verbatim}
5
7
11
101
131
151
181
191
313
353
373
383
\end{verbatim}

\paragraph{SAMPLE INPUT (file sprime.in)}
\begin{verbatim}
4
\end{verbatim}

\paragraph{SAMPLE OUTPUT (file sprime.out)}
\begin{verbatim}
2333
2339
2393
2399
2939
3119
3137
3733
3739
3793
3797
5939
7193
7331
7333
7393
\end{verbatim}

\paragraph{SAMPLE INPUT(checker.in)}
\begin{verbatim}
6
\end{verbatim}

\paragraph{SAMPLE OUTPUT(checker.out)}
\begin{verbatim}
2 4 6 1 3 5
3 6 2 5 1 4
4 1 5 2 6 3
\end{verbatim}

\paragraph{SAMPLE INPUT (file castle.in)}
\begin{verbatim}
7 4
11 6 11 6 3 10 6
7 9 6 13 5 15 5
1 10 12 7 13 7 5
13 11 10 8 10 12 13
\end{verbatim}

\paragraph{SAMPLE OUTPUT (file castle.out)}
\begin{verbatim}
5
9
16
4 1 E
\end{verbatim}

\paragraph{SAMPLE INPUT (file frac1.in)}
\begin{verbatim}
5
\end{verbatim}

\paragraph{SAMPLE OUTPUT (file frac1.out)}
\begin{verbatim}
0/1
1/5
1/4
1/3
2/5
1/2
3/5
2/3
3/4
4/5
1/1
\end{verbatim}

\paragraph{SAMPLE INPUT (file sort3.in)}
\begin{verbatim}
9
\end{verbatim}

\paragraph{SAMPLE OUTPUT (file sort3.out)}
\begin{verbatim}
4
\end{verbatim}

\paragraph{SAMPLE INPUT (file holstein.in)}
\begin{verbatim}
4
100 200 300 400
3
50 50 50 50
200 300 200 300
900 150 389 399
\end{verbatim}

\paragraph{SAMPLE OUTPUT (file holstein.out)}
\begin{verbatim}
2 1 3
\end{verbatim}

\paragraph{SAMPLE INPUT (file hamming.in)}
\begin{verbatim}
16 7 3
\end{verbatim}

\paragraph{SAMPLE OUTPUT (file hamming.out)}
\begin{verbatim}
0 7 25 30 42 45 51 52 75 76
82 85 97 102 120 127
\end{verbatim}

\paragraph{SAMPLE INPUT(preface.in)}
\begin{verbatim}
5
\end{verbatim}

\paragraph{SAMPLE OUTPUT(preface.out)}
\begin{verbatim}
I 7
V 2
\end{verbatim}

\paragraph{SAMPLE INPUT (file subset.in)}
\begin{verbatim}
7
\end{verbatim}

\paragraph{SAMPLE OUTPUT (file subset.out)}
\begin{verbatim}
4
\end{verbatim}

\paragraph{SAMPLE INPUT (file runround.in)}
\begin{verbatim}
81361
\end{verbatim}

\paragraph{SAMPLE OUTPUT (file runround.out)}
\begin{verbatim}
81362
\end{verbatim}

\paragraph{SAMPLE INPUT (file lamps.in)}
\begin{verbatim}
10
1
-1
7 -1
\end{verbatim}

\paragraph{SAMPLE OUTPUT (file lamps.out)}
\begin{verbatim}
0000000000
0101010101
0110110110
\end{verbatim}

\paragraph{SAMPLE INPUT (file prefix.in)}
\begin{verbatim}
A AB BA CA BBC
ABABACABAABC
\end{verbatim}

\paragraph{SAMPLE OUTPUT (file prefix.out)}
\begin{verbatim}
11
\end{verbatim}

\paragraph{SAMPLE INPUT (file nocows.in)}
\begin{verbatim}
5 3
\end{verbatim}

\paragraph{SAMPLE OUTPUT (file nocows.out)}
\begin{verbatim}
2
\end{verbatim}

\paragraph{SAMPLE INPUT (file zerosum.in)}
\begin{verbatim}
7
\end{verbatim}

\paragraph{SAMPLE OUTPUT (file zerosum.out)}
\begin{verbatim}
1+2-3+4-5-6+7
1+2-3-4+5+6-7
1-2 3+4+5+6+7
1-2 3-4 5+6 7
1-2+3+4-5+6-7
1-2-3-4-5+6+7
\end{verbatim}

\paragraph{SAMPLE INPUT (file money.in)}
\begin{verbatim}
3 10
1 2 5
\end{verbatim}

\paragraph{SAMPLE OUTPUT (file money.out)}
\begin{verbatim}
10
\end{verbatim}

\paragraph{SAMPLE INPUT (file concom.in)}
\begin{verbatim}
3
1 2 80
2 3 80
3 1 20
\end{verbatim}

\paragraph{SAMPLE OUTPUT (file concom.out)}
\begin{verbatim}
1 2
1 3
2 3
\end{verbatim}

\paragraph{SAMPLE INPUT (file ttwo.in)}
\begin{verbatim}
*...*.....
......*...
...*...*..
..........
...*.F....
*.....*...
...*......
..C......*
...*.*....
.*.*......
\end{verbatim}

\paragraph{SAMPLE OUTPUT (file ttwo.out)}
\begin{verbatim}
49
\end{verbatim}

\paragraph{SAMPLE INPUT (file maze1.in)}
\begin{verbatim}
5 3
+-+-+-+-+-+
|
|
+-+ +-+ + +
|
| | |
+ +-+-+ + +
| |
|
+-+ +-+-+-+
\end{verbatim}

\paragraph{SAMPLE OUTPUT (file maze1.out)}
\begin{verbatim}
9
\end{verbatim}

\paragraph{SAMPLE INPUT (file comehome.in)}
\begin{verbatim}
5
A d 6
B d 3
C e 9
d Z 8
e Z 3
\end{verbatim}

\paragraph{SAMPLE OUTPUT (file comehome.out)}
\begin{verbatim}
B 11
\end{verbatim}

\paragraph{SAMPLE INPUT (file cowtour.in)}
\begin{verbatim}
8
10 10
15 10
20 10
15 15
20 15
30 15
25 10
30 10
01000000
10111000
01001000
01001000
01110000
00000010
00000101
\end{verbatim}

\paragraph{SAMPLE OUTPUT (file cowtour.out)}
\begin{verbatim}
22.071068
\end{verbatim}

\paragraph{SAMPLE INPUT (file fracdec.in)}
\begin{verbatim}
45 56
\end{verbatim}

\paragraph{SAMPLE OUTPUT (file fracdec.out)}
\begin{verbatim}
0.803(571428)
\end{verbatim}

\paragraph{SAMPLE INPUT (file agrinet.in)}
\begin{verbatim}
4
0 4 9 21
4 0 8 17
9 8 0 16
21 17 16 0
\end{verbatim}

\paragraph{SAMPLE OUTPUT (file agrinet.out)}
\begin{verbatim}
28
\end{verbatim}

\paragraph{SAMPLE INPUT (file inflate.in)}
\begin{verbatim}
300 4
100 60
250 120
120 100
35 20
\end{verbatim}

\paragraph{SAMPLE OUTPUT (file inflate.out)}
\begin{verbatim}
605
\end{verbatim}

\paragraph{SAMPLE INPUT (file humble.in)}
\begin{verbatim}
4 19
2 3 5 7
\end{verbatim}

\paragraph{SAMPLE OUTPUT (file humble.out)}
\begin{verbatim}
27
\end{verbatim}

\paragraph{SAMPLE INPUT (file rect1.in)}
\begin{verbatim}
20 20 3
2 2 18 18 2
0 8 19 19 3
8 0 10 19 4
\end{verbatim}

\paragraph{SAMPLE OUTPUT (file rect1.out)}
\begin{verbatim}
1 91
2 84
3 187
4 38
\end{verbatim}

\paragraph{SAMPLE INPUT (file contact.in)}
\begin{verbatim}
2 4 10
01010010010001000111101100001010011001111000010010011110010000000
\end{verbatim}

\paragraph{SAMPLE OUTPUT (file contact.out)}
\begin{verbatim}
23
00
15
01 10
12
100
11
11 000 001
10
010
8
0100
7
0010 1001
6
111 0000
5
011 110 1000
4
0001 0011 1100
\end{verbatim}

\paragraph{SAMPLE INPUT (file stamps.in)}
\begin{verbatim}
5 2
1 3
\end{verbatim}

\paragraph{SAMPLE OUTPUT (file stamps.out)}
\begin{verbatim}
13
\end{verbatim}

\paragraph{SAMPLE INPUT (file fact4.in)}
\begin{verbatim}
7
\end{verbatim}

\paragraph{SAMPLE OUTPUT (file fact4.out)}
\begin{verbatim}
4
\end{verbatim}

\paragraph{SAMPLE INPUT (file kimbits.in)}
\begin{verbatim}
5
\end{verbatim}

\paragraph{SAMPLE OUTPUT (file kimbits.out)}
\begin{verbatim}
10011
\end{verbatim}

\paragraph{SAMPLE INPUT (file spin.in)}
\begin{verbatim}
30 1 0 120
50 1 150 90
60 1 60 90
70 1 180 180
90 1 180 60
\end{verbatim}

\paragraph{SAMPLE OUTPUT (file spin.out)}
\begin{verbatim}
9
\end{verbatim}

\paragraph{SAMPLE INPUT (file ratios.in)}
\begin{verbatim}
3 4 5 1 2 3 3 7 1 2 1 2
\end{verbatim}

\paragraph{SAMPLE OUTPUT (file ratios.out)}
\begin{verbatim}
8 1 5 7
\end{verbatim}

\paragraph{SAMPLE INPUT (file msquare.in)}
\begin{verbatim}
2 6 8 4 5 7 3 1
\end{verbatim}

\paragraph{SAMPLE OUTPUT (file msquare.out)}
\begin{verbatim}
7 BCABCCB
\end{verbatim}

\paragraph{SAMPLE INPUT (file butter.in)}
\begin{verbatim}
3 4 5
2
3
4
1 2 1
1 3 5
2 3 7
2 4 3
3 4 5
P2
P1 @--1--@ C1
\
|\
\
| \
5 7 3
\ |
\
\|
\ C3
C2 @--5--@
P3
P4
}
\end{verbatim}

\paragraph{SAMPLE OUTPUT (file butter.out)}
\begin{verbatim}
8
\end{verbatim}

\paragraph{SAMPLE INPUT(fence.in)}
\begin{verbatim}
9
1 2
2 3
3 4
4 2
4 5
2 5
5 6
5 7
4 6
\end{verbatim}

\paragraph{SAMPLE OUTPUT(fence.out)}
\begin{verbatim}
1
\end{verbatim}

\paragraph{SAMPLE INPUT (file shopping.in)}
\begin{verbatim}
2
1 7 3 5
2 7 1 8 2 10
2
7 3 2
8 2 5
\end{verbatim}

\paragraph{SAMPLE OUTPUT (file shopping.out)}
\begin{verbatim}
14
\end{verbatim}

\paragraph{SAMPLE INPUT (file camelot.in)}
\begin{verbatim}
8 8
D 4
A 3 A 8
H 1 H 8
\end{verbatim}

\paragraph{SAMPLE OUTPUT (file camelot.out)}
\begin{verbatim}
10
\end{verbatim}

\paragraph{SAMPLE INPUT (file range.in)}
\begin{verbatim}
6
101111
001111
111111
001111
101101
111001
\end{verbatim}

\paragraph{SAMPLE OUTPUT (file range.out)}
\begin{verbatim}
2 10
3 4
4 1
\end{verbatim}

\paragraph{SAMPLE INPUT (file game1.in)}
\begin{verbatim}
6
4 7 2 9
5 2
\end{verbatim}

\paragraph{SAMPLE OUTPUT (file game1.out)}
\begin{verbatim}
18 11
\end{verbatim}

\paragraph{SAMPLE INPUT (file fence4.in)}
\begin{verbatim}
13
5 5
0 0
7 0
5 2
7 5
5 7
3 5
4 9
1 8
2 5
0 9
-2 7
0 3
-3 1
\end{verbatim}

\paragraph{SAMPLE OUTPUT (file fence4.out)}
\begin{verbatim}
7
0 0 7 0
5 2 7 5
7 5 5 7
5 7 3 5
-2 7 0 3
0 0 -3 1
0 3 -3 1
\end{verbatim}

\paragraph{SAMPLE INPUT (file heritage.in)}
\begin{verbatim}
ABEDFCHG
CBADEFGH
D
/ \
E
\end{verbatim}

\paragraph{SAMPLE OUTPUT (file heritage.out)}
\begin{verbatim}
AEFDBHGC
\end{verbatim}

\paragraph{SAMPLE INPUT (file fence3.in)}
\begin{verbatim}
3
0 0 0 1
2 0 2 1
0 3 2 3
\end{verbatim}

\paragraph{SAMPLE OUTPUT (file fence3.out)}
\begin{verbatim}
1.6 3.7
\end{verbatim}

\paragraph{SAMPLE INPUT (file rockers.in)}
\begin{verbatim}
4 5 2
4 3 4 2
\end{verbatim}

\paragraph{SAMPLE OUTPUT (file rockers.out)}
\begin{verbatim}
3
\end{verbatim}

\paragraph{SAMPLE INPUT (file nuggets.in)}
\begin{verbatim}
3
\end{verbatim}

\paragraph{SAMPLE OUTPUT (file nuggets.out)}
\begin{verbatim}
17
\end{verbatim}

\paragraph{SAMPLE INPUT (file fence8.in)}
\begin{verbatim}
4
\end{verbatim}

\paragraph{SAMPLE OUTPUT (file fence8.out)}
\begin{verbatim}
7
\end{verbatim}

\paragraph{SAMPLE INPUT (file fence6.in)}
\begin{verbatim}
10
1 16 2 2
2 7
10 6
2 3 2 2
1 7
8 3
3 3 2 1
8 2
4
4 8 1 3
3
9 10 5
5 8 3 1
9 10 4
6
6 6 1 2
5
1 10
7 5 2 2
1 2
8 9
8 4 2 2
2 3
7 9
9 5 2 3
7 8
4 5 10
10 10 2 3
1 6
4 9 5
\end{verbatim}

\paragraph{SAMPLE OUTPUT (file fence6.out)}
\begin{verbatim}
12
\end{verbatim}

\paragraph{SAMPLE INPUT(file cryptcow.in)}
\begin{verbatim}
Begin the EscCution at the BreOape execWak of Dawn
\end{verbatim}

\paragraph{SAMPLE OUTPUT(file cryptcow.out)}
\begin{verbatim}
1 1
\end{verbatim}

\paragraph{SAMPLE INPUT (file ditch.in)}
\begin{verbatim}
5 4
1 2 40
1 4 20
2 4 20
2 3 30
3 4 10
\end{verbatim}

\paragraph{SAMPLE OUTPUT (file ditch.out)}
\begin{verbatim}
50
\end{verbatim}

\paragraph{SAMPLE INPUT (file stall4.in)}
\begin{verbatim}
5 5
2 2 5
3 2 3 4
2 1 5
3 1 2 5
1 2
\end{verbatim}

\paragraph{SAMPLE OUTPUT (file stall4.out)}
\begin{verbatim}
4
\end{verbatim}

\paragraph{SAMPLE INPUT (file job.in)}
\begin{verbatim}
5 2 3
1 1 3 1 4
\end{verbatim}

\paragraph{SAMPLE OUTPUT (file job.out)}
\begin{verbatim}
3 5
\end{verbatim}

\paragraph{SAMPLE INPUT (file cowcycle.in)}
\begin{verbatim}
2 5
39 62 12 28
\end{verbatim}

\paragraph{SAMPLE OUTPUT (file cowcycle.out)}
\begin{verbatim}
57
39 53
12 13 15 23 27
Comment
2 5
39 62 12 28
39/12 = 3.25000000000000000000
39/13 = 3.00000000000000000000
39/14 = 2.78571428571428571428
39/15 = 2.60000000000000000000
39/16 = 2.43750000000000000000
40/12 = 3.33333333333333333333
40/13 = 3.07692307692307692307
40/14 = 2.85714285714285714285
40/15 = 2.66666666666666666666
40/16 = 2.50000000000000000000
39/16 = 2.43750000000000000000
40/16 = 2.50000000000000000000
39/15 = 2.60000000000000000000
40/15 = 2.66666666666666666666
39/14 = 2.78571428571428571428
40/14 = 2.85714285714285714285
39/13 = 3.00000000000000000000
40/13 = 3.07692307692307692307
39/12 = 3.25000000000000000000
40/12 = 3.33333333333333333333
2.43750000000000000000 - 2.50000000000000000000 = 0.06250000000000000000
2.50000000000000000000 - 2.60000000000000000000 = 0.10000000000000000000
2.60000000000000000000 - 2.66666666666666666666 = 0.06666666666666666666
2.66666666666666666666 - 2.78571428571428571428 = 0.11904761904761904762
2.78571428571428571428 - 2.85714285714285714285 = 0.07142857142857142857
2.85714285714285714285 - 3.00000000000000000000 = 0.14285714285714285715
3.00000000000000000000 - 3.07692307692307692307 = 0.07692307692307692307
3.07692307692307692307 - 3.25000000000000000000 = 0.17307692307692307693
3.25000000000000000000 - 3.33333333333333333333 = 0.08333333333333333333
\end{verbatim}

\paragraph{SAMPLE INPUT (file buylow.in)}
\begin{verbatim}
12
68 69 54 64 68 64 70 67
78 62 98 87
\end{verbatim}

\paragraph{SAMPLE OUTPUT (file buylow.out)}
\begin{verbatim}
4 2
\end{verbatim}

\paragraph{SAMPLE INPUT (file prime3.in)}
\begin{verbatim}
11 1
\end{verbatim}

\paragraph{SAMPLE OUTPUT (file prime3.out)}
\begin{verbatim}
11351
14033
30323
53201
13313
11351
33203
30323
14033
33311
13313
13043
32303
50231
13331
\end{verbatim}

\paragraph{SAMPLE INPUT (file race3.in)}
\begin{verbatim}
1 2 -2
3 -2
3 -2
5 4 -2
6 4 -2
6 -2
7 8 -2
9 -2
5 9 -2
-2
-1
\end{verbatim}

\paragraph{SAMPLE OUTPUT (file race3.out)}
\begin{verbatim}
2 3 6
1 3
\end{verbatim}

\paragraph{SAMPLE INPUT (file lgame.in)}
\begin{verbatim}
prmgroa
\end{verbatim}

\paragraph{SAMPLE INPUT (file lgame.dict)}
\begin{verbatim}
profile
program
prom
rag
ram
rom
.
\end{verbatim}

\paragraph{SAMPLE OUTPUT (file lagame.out)}
\begin{verbatim}
24
program
prom rag
\end{verbatim}

\paragraph{SAMPLE INPUT (file shuttle.in)}
\begin{verbatim}
3
\end{verbatim}

\paragraph{SAMPLE OUTPUT (file shuttle.out)}
\begin{verbatim}
3 5 6 4 2 1 3 5 7 6 4 2 3 5 4
\end{verbatim}

\paragraph{SAMPLE INPUT (file milk6.in)}
\begin{verbatim}
4 5
1 3 100
3 2 50
2 4 60
1 2 40
2 3 80
\end{verbatim}

\paragraph{SAMPLE OUTPUT (file milk6.out)}
\begin{verbatim}
60 1 3
\end{verbatim}

\paragraph{SAMPLE INPUT (file frameup.in)}
\begin{verbatim}
9 8
CCC....
ECBCBB..
DCBCDB..
DCCC.B..
D.B.ABAA
D.BBBB.A
DDDDAD.A
E...AAAA
EEEEEE..
\end{verbatim}

\paragraph{SAMPLE OUTPUT (file frameup.out)}
\begin{verbatim}
EDABC
\end{verbatim}

\paragraph{SAMPLE INPUT (file fc.in)}
\begin{verbatim}
4
4 8
4 12
5 9.3
7 8
\end{verbatim}

\paragraph{SAMPLE OUTPUT (file fc.out)}
\begin{verbatim}
12.00
\end{verbatim}

\paragraph{SAMPLE INPUT (file starry.in)}
\begin{verbatim}
23
15
10001000000000010000000
01111100011111000101101
01000000010001000111111
00000000010101000101111
00000111010001000000000
00001001011111000000000
10000001000000000000000
00101000000111110010000
00001000000100010011111
00000001110101010100010
00000100110100010000000
00010001110111110000000
00100001110000000100000
00001000100001000100101
00000001110001000111000
\end{verbatim}

\paragraph{SAMPLE OUTPUT (file starry.out)}
\begin{verbatim}
a000a0000000000b0000000
0aaaaa000ccccc000d0dd0d
0a0000000c000c000dddddd
000000000c0b0c000d0dddd
00000eee0c000c000000000
0000e00e0ccccc000000000
b000000e000000000000000
00b0f000000ccccc00a0000
0000f000000c000c00aaaaa
0000000ddd0c0b0c0a000a0
00000b00dd0c000c0000000
000g000ddd0ccccc0000000
00g0000ddd0000000e00000
0000b000d0000f000e00e0b
0000000ddd000f000eee000
Constraints
\end{verbatim}

\paragraph{SAMPLE INPUT (file theme.in)}
\begin{verbatim}
30
25 27 30 34 39 45 52 60 69 79 69 60 52 45 39 34 30 26 22 18
82 78 74 70 66 67 64 60 65 80
\end{verbatim}

\paragraph{SAMPLE OUTPUT (file theme.out)}
\begin{verbatim}
5
\end{verbatim}

\paragraph{SAMPLE INPUT (file snail.in)}
\begin{verbatim}
8 4 E2 A6 G1 F5
\end{verbatim}

\paragraph{SAMPLE OUTPUT (file snail.out)}
\begin{verbatim}
33
\end{verbatim}

\paragraph{SAMPLE INPUT (file fence3.in)}
\begin{verbatim}
3
0 0 0 1
2 0 2 1
0 3 2 3
\end{verbatim}

\paragraph{SAMPLE OUTPUT (file fence3.out)}
\begin{verbatim}
1.6 3.7
\end{verbatim}

\paragraph{SAMPLE INPUT (file wissqu.in)}
\begin{verbatim}
ABAC
DCDE
BEBC
CADE
\end{verbatim}

\paragraph{SAMPLE OUTPUT (file wissqu.out)}
\begin{verbatim}
D 4 1
C 4 2
A 3 1
A 3 3
B 2 4
B 3 2
B 4 4
E 2 1
E 2 3
D 1 4
D 2 2
C 1 1
C 1 3
A 1 2
E 4 3
D 3 4
14925
\end{verbatim}

\paragraph{SAMPLE INPUT (file milk4.in)}
\begin{verbatim}
16 3 3 5 7
\end{verbatim}

\paragraph{SAMPLE OUTPUT (file milk4.out)}
\begin{verbatim}
2 3 5
\end{verbatim}

\paragraph{SAMPLE INPUT (file window.in)}
\begin{verbatim}
w(a,10,132,20,12)
w(b,8,76,124,15)
s(a)
\end{verbatim}

\paragraph{SAMPLE OUTPUT (file window.out)}
\begin{verbatim}
49.167
\end{verbatim}

\paragraph{SAMPLE INPUT (file schlnet.in)}
\begin{verbatim}
5
2 4 3 0
4 5 0
0
0
1 0
\end{verbatim}

\paragraph{SAMPLE OUTPUT (file schlnet.out)}
\begin{verbatim}
1
\end{verbatim}

\paragraph{SAMPLE INPUT (file bigbrn.in)}
\begin{verbatim}
8 3
2 2
2 6
6 3
\end{verbatim}

\paragraph{SAMPLE OUTPUT (file bigbrn.out)}
\begin{verbatim}
5
\end{verbatim}

\paragraph{SAMPLE INPUT (file latin.in)}
\begin{verbatim}
5
\end{verbatim}

\paragraph{SAMPLE OUTPUT (file latin.out)}
\begin{verbatim}
1344
\end{verbatim}

\paragraph{SAMPLE INPUT (file tour.in)}
\begin{verbatim}
8 9
Vancouver
Yellowknife
Edmonton
Calgary
Winnipeg
Toronto
Montreal
Halifax
Vancouver Edmonton
Vancouver Calgary
Calgary Winnipeg
Winnipeg Toronto
Toronto Halifax
Montreal Halifax
Edmonton Montreal
Edmonton Yellowknife
Edmonton Calgary
\end{verbatim}

\paragraph{SAMPLE OUTPUT (file tour.out)}
\begin{verbatim}
7
\end{verbatim}

\paragraph{SAMPLE INPUT (file char.in)}
\begin{verbatim}
font.in
540
00000000000000000000
00000000000000000000
00000000000000000000
00000000000000000000
00000000000000000000
00000000000000000000
00000000000000000000
00000000000000000000
00000000000000000000
00000000000000000000
00000000000000000000
00000000000000000000
00000000000000000000
00000000000000000000
00000000000000000000
00000000000000000000
00000000000000000000
00000000000000000000
00000000000000000000
00000000000000000000
00000000000000000000
00000000000000000000
00000000000000000000
00000011100000000000
“a”
char.in
19
00000000000000000000
00000000000000000000
00000000000000000000
00000011100000000000
00100111011011000000
00001111111001100000
00001110001100100000
00001100001100010000
00001100000100010000
00000100000100010000
00000010000000110000
00001111011111110000
00001111111111110000
00001111111111000000
00001000010000000000
00000000000000000000
00000000000001000000
00000000000000000000
00000000000000000000
00000111111011000000
00001111111001100000
00001110001100100000
00001100001100010000
00001100000100010000
00000100000100010000
00000010000000110000
00000001000001110000
00001111111111110000
00001111111111110000
00001111111111000000
00001000000000000000
00000000000000000000
00000000000000000000
00000000000000000000
00000000000000000000
\end{verbatim}

\paragraph{SAMPLE OUTPUT (file char.out)}
\begin{verbatim}
a
\end{verbatim}

\paragraph{SAMPLE INPUT (file betsy.in)}
\begin{verbatim}
3
\end{verbatim}

\paragraph{SAMPLE OUTPUT (file betsy.out)}
\begin{verbatim}
2
\end{verbatim}

\paragraph{SAMPLE INPUT (file telecow.in)}
\begin{verbatim}
3 2 1 2
1 3
2 3
\end{verbatim}

\paragraph{SAMPLE OUTPUT(file telecow.out)}
\begin{verbatim}
1
\end{verbatim}

\paragraph{SAMPLE INPUT (file picture.in)}
\begin{verbatim}
7 -15 0 5 10 -5 8 20 25 15 -4 24 14 0 -6 16 4 2 15 10 22 30 10 36 20 34 0 40 16
\end{verbatim}

\paragraph{SAMPLE OUTPUT (file picture.out)}
\begin{verbatim}
228
\end{verbatim}

\paragraph{SAMPLE INPUT (file hidden.in)}
\begin{verbatim}
7
alabala
\end{verbatim}

\paragraph{SAMPLE OUTPUT (file hidden.out)}
\begin{verbatim}
6
\end{verbatim}

\paragraph{SAMPLE INPUT \#1(file twofive.in)}
\begin{verbatim}
N
\end{verbatim}

\paragraph{SAMPLE INPUT \#2}
\begin{verbatim}
W
ABCDEFGHIJKLMNOPQRSUTVWXY
\end{verbatim}

\paragraph{SAMPLE OUTPUT \#1(file twofive.out)}
\begin{verbatim}
ABCDEFGHIJKLMNOPQRSUTVWXY
\end{verbatim}

\paragraph{SAMPLE OUTPUT \#2}
\begin{verbatim}
2
\end{verbatim}

\paragraph{SAMPLE INPUT (file vans.in)}
\begin{verbatim}
4
\end{verbatim}

\paragraph{SAMPLE OUTPUT (file vans.out)}
\begin{verbatim}
12
\end{verbatim}

\paragraph{SAMPLE INPUT (file rectbarn.in)}
\begin{verbatim}
3 4 2
1 3
2 1
\end{verbatim}

\paragraph{SAMPLE OUTPUT (file rectbarn.out)}
\begin{verbatim}
6
\end{verbatim}

\paragraph{SAMPLE INPUT (file cowxor.in)}
\begin{verbatim}
5
1
0
\end{verbatim}

\paragraph{SAMPLE OUTPUT (file cowxor.out)}
\begin{verbatim}
6 4 5
\end{verbatim}

\end{document}
